% IEEE include snippet for Table D. Scope and Non-Claims
% Requires in preamble: \usepackage{booktabs,tabularx,array}
\begin{table}[t]
\centering
\footnotesize
\setlength{\tabcolsep}{3.4pt}
\renewcommand{\arraystretch}{1.08}
\caption{Explicit scope boundary for the current RAASA evidence package.}
\label{tab:scope-nonclaims}
\begin{tabularx}{\columnwidth}{p{0.22\columnwidth} c >{\raggedright\arraybackslash}X p{0.21\columnwidth}}
\toprule
Claim area & Supported? & Evidence basis & Note \\
\midrule
Adaptive vs static local trade-off & Yes & Local Docker baseline comparison & Core thesis anchor \\
Fresh-account cloud reproducibility & Yes & Single-node replay on a clean AWS Free-plan account & Bounded replay only \\
Bounded multi-node K3s continuity & Yes & 3-node soak, repeated matrix, degraded-mode handling, and drain/reschedule & K3s only; bounded envelope \\
Managed Kubernetes / EKS robustness & No & No EKS evidence bundle exists & Must remain a non-claim \\
Multi-tenant safety & No & No multi-tenant evaluation exists & Must remain a non-claim \\
Production readiness & No & Bounded research-prototype evidence only & Explicitly avoid this claim \\
Broad exfiltration-prevention guarantee & No & Only bounded adversarial matrix coverage exists & Specific workloads only \\
High availability / fault tolerance & No & Drain/reschedule continuity exists, but no HA or SLA evidence & Do not overread the worker-drain result \\
\bottomrule
\end{tabularx}
\end{table}

% Notes
% - This table is helpful in an appendix, rebuttal package, or artifact companion.
% - It prevents readers from inferring claims larger than the evidence base.
